\section*{Capítulo \ref{ch:g10:exercise-and-energy} --- O exercício e as necessidades energéticas do corpo}

\begin{solution}{exo:g10:exercise-and-energy:1}
A energia gasta em repouso completo para manter o corpo vivo ---
coração, cérebro, temperatura, manutenção das \omterm{def:g10:cells-common-unit:cell}{células}: cerca de
\qty{80}{W}, uns \qty{7000}{kJ} por dia.
\end{solution}

\begin{solution}{exo:g10:exercise-and-energy:2}
$2.0 \times 20 = \qty{40}{kJ}$ por minuto, ou seja $40000/60 \approx
\qty{670}{W}$.
\end{solution}

\begin{solution}{exo:g10:exercise-and-energy:3}
A reserva de fosfato do músculo (cerca de dez segundos); a \omterm{def:g10:cell-metabolism:fermentation}{fermentação}
láctica (um a dois minutos de esforço máximo); a \omterm{prop:g10:cell-metabolism:respiration}{respiração celular}
(horas).
\end{solution}

\begin{solution}{exo:g10:exercise-and-energy:4}
A maior taxa com que o corpo consegue captar e usar oxigênio.
$3500/70 = \qty{50}{mL/min}$ por quilograma.
\end{solution}

\begin{solution}{exo:g10:exercise-and-energy:5}
Nos músculos, cerca de \qty{400}{g}, abastecendo o músculo que o
armazena; no fígado, cerca de \qty{100}{g}, liberado como glicose no
sangue para o cérebro e os demais tecidos.
\end{solution}

\begin{solution}{exo:g10:exercise-and-energy:6}
\qty{2.7}{L/min}: $2.7 \times 20 = \qty{54}{kJ}$ por minuto, ou seja
\qty{900}{W}. Calor: $900 - 200 = \qty{700}{W}$.
\end{solution}

\begin{solution}{exo:g10:exercise-and-energy:7}
Trabalho $mgh = 60 \times 9.8 \times 800 \approx \qty{470}{kJ}$. Com
25\% de rendimento, os músculos gastaram cerca de \qty{1900}{kJ}. Ao
longo de \qty{7200}{s}: cerca de \qty{260}{W} acima do repouso.
\end{solution}

\begin{solution}{exo:g10:exercise-and-energy:8}
Um tiro de 100 metros funciona quase inteiramente com a reserva de
fosfato: nenhuma dívida de oxigênio, pouco ácido láctico. Uma prova de
400 metros dura cerca de um minuto e tira metade da energia da
\omterm{def:g10:cell-metabolism:fermentation}{fermentação} láctica: o ácido láctico se acumula (a dor) e a respiração
que precisa eliminá-lo mantém o corredor ofegante por minutos depois.
\end{solution}

\begin{solution}{exo:g10:exercise-and-energy:9}
Energia de \omterm{def:g10:chemistry-of-life:families}{carboidrato} $0.6 \times 3000 = \qty{1800}{kJ}$; \omterm{def:g10:exercise-and-energy:glycogen}{glicogênio}
$1800/17 \approx \qty{106}{g}$; água liberada, cerca de \qty{320}{g}.
\end{solution}

\begin{solution}{exo:g10:exercise-and-energy:10}
Potência metabólica de cerca de \qty{600}{W}: $\num{1.1e6}/600 \approx
\qty{1830}{s}$, cerca de 30 minutos.
\end{solution}

\begin{solution}{exo:g10:exercise-and-energy:11}
Três quartos da energia liberada nos músculos são calor, várias
centenas de watts durante um exercício puxado. A transpiração o
elimina: evaporar água consome cerca de \qty{2.4}{kJ} por grama, e o
fluxo de sangue para a pele leva o calor até lá. Sem transpiração, a
temperatura do corpo subiria um grau a cada poucos minutos.
\end{solution}

\begin{solution}{exo:g10:exercise-and-energy:12}
Na subida, a potência necessária é proporcional à massa corporal: o
sujeito de \qty{60}{kg} tem \qty{60}{mL/min} por quilograma contra
\num{40}, e sobe mais rápido. Na bicicleta de laboratório a potência
não depende da massa, os dois têm os mesmos \qty{3.6}{L/min}, e os dois
conseguem sustentar os mesmos watts.
\end{solution}

\begin{solution}{exo:g10:exercise-and-energy:13}
Os \qty{0.4}{L/min} de equivalente em oxigênio que faltam são
fornecidos pela \omterm{def:g10:cell-metabolism:fermentation}{fermentação} láctica nos músculos. O ácido láctico se
acumula; em poucos minutos os músculos ardem e o corredor é obrigado a
baixar para um ritmo que a respiração sozinha consiga cobrir.
\end{solution}

\begin{solution}{exo:g10:exercise-and-energy:14}
Os tiros curtos e máximos funcionam com a reserva de fosfato e a
\omterm{def:g10:cell-metabolism:fermentation}{fermentação}, que não usam gordura nenhuma; a gordura só é queimada pela
respiração, que leva minutos para chegar ao ritmo máximo e só domina no
exercício moderado e prolongado. A energia total de alguns tiros
também é pequena comparada com a de uma hora de esforço constante.
\end{solution}

\begin{solution}{exo:g10:exercise-and-energy:15}
Déficit $25000 - 8500 - 6000 = \qty{10500}{kJ}$, ou seja cerca de
\qty{280}{g} de gordura. Comer durante a etapa mantém a glicose
sanguínea alta para o cérebro e poupa o \omterm{def:g10:exercise-and-energy:glycogen}{glicogênio} para as subidas,
onde a potência necessária está muito acima do que a gordura sozinha
consegue fornecer; a gordura cobre a parte constante do dia.
\end{solution}

\begin{solution}{pb:g10:exercise-and-energy:1}
\textbf{1.} $4.0 \times 60 \times 42.2 \approx \qty{10100}{kJ}$.

\textbf{2.} $\qty{3}{h}\,\qty{30}{min} = \qty{12600}{s}$:
$\num{1.01e7}/12600 \approx \qty{800}{W}$.

\textbf{3.} $10100/20 \approx \qty{506}{L}$; em \qty{210}{min}, cerca
de \qty{2.4}{L/min}.

\textbf{4.} $2400/60 = \qty{40}{mL/min}$ por quilograma, ou seja $40/52
\approx 77\%$ do máximo dela.

\textbf{5.} $(2.4 - 0.2)/2.4 \approx 92\%$.

\textbf{6.} $380 \times 17 = \qty{6460}{kJ}$.

\textbf{7.} A prova custa \qty{240}{kJ} por quilômetro: $6460/240
\approx \qty{27}{km}$.

\textbf{8.} Uso de \omterm{def:g10:chemistry-of-life:families}{carboidrato} $0.8 \times 240 = \qty{192}{kJ/km}$:
$6460/192 \approx \qty{34}{km}$ --- o muro.

\textbf{9.} Energia da gordura $0.2 \times 10100 \approx \qty{2020}{kJ}$,
ou seja $2020/38 \approx \qty{53}{g}$: cerca de 0.6\% dos \qty{9}{kg}
dela.

\textbf{10.} A gordura só é queimada pela respiração e a uma taxa
limitada: ela consegue fornecer aproximadamente metade da potência que
o ritmo dela exige. Sem \omterm{def:g10:exercise-and-energy:glycogen}{glicogênio} os músculos não conseguem fabricar
ATP depressa o bastante, e o ritmo cai.

\textbf{11.} $60 \times 3.5 = \qty{210}{g}$, ou seja $210 \times 17 \approx
\qty{3570}{kJ}$, cobrindo $3570/192 \approx \qty{19}{km}$ de uso de
\omterm{def:g10:chemistry-of-life:families}{carboidrato}.

\textbf{12.} \omterm{def:g10:chemistry-of-life:families}{Carboidrato} disponível $6460 + 3570 = \qty{10030}{kJ}$:
$10030/192 \approx \qty{52}{km}$, além da chegada. Sim.

\textbf{13.} $+\qty{200}{g}$ de \omterm{def:g10:exercise-and-energy:glycogen}{glicogênio} $+ \qty{600}{g}$ de água:
\qty{0.8}{kg}, custando $0.8 \times 4.0 \approx \qty{3}{kJ}$ a mais por
quilômetro --- desprezível diante de 240.

\textbf{14.} \omterm{def:g10:exercise-and-energy:glycogen}{Glicogênio} $580 \times 17 = \qty{9860}{kJ}$ mais a bebida:
cerca de \qty{13400}{kJ}; uso de \omterm{def:g10:chemistry-of-life:families}{carboidrato} $0.55 \times 240 =
\qty{132}{kJ/km}$: mais de \qty{100}{km}. O muro se deslocou para bem
além da chegada, com boa margem.

\textbf{15.} O cérebro só queima glicose, fornecida pelo \omterm{def:g10:exercise-and-energy:glycogen}{glicogênio} do
fígado. Quando essa reserva se esvazia, a glicose sanguínea cai, e o
cérebro sinaliza a emergência como tontura e cansaço avassalador, seja
o que for que os músculos ainda guardem.

\textbf{16.} $3.6 \times 55 \times 42.2 \approx \qty{8360}{kJ}$;
$8360/20 = \qty{418}{L}$ em \qty{125}{min}: cerca de \qty{3.3}{L/min}.

\textbf{17.} $3340/55 \approx \qty{61}{mL/min}$ por quilograma, ou seja
$61/80 \approx 76\%$ --- a mesma fração de Nádia. O teto é que difere,
não a fração dele que é usada.

\textbf{18.} A economia é o custo por quilograma por quilômetro. Para um
corredor de \qty{55}{kg}, $(4.0 - 3.6) \times 55 \times 42.2 \approx
\qty{930}{kJ}$ economizados na prova, cerca de 9\%.

\textbf{19.} \omterm{def:g10:exercise-and-energy:glycogen}{Glicogênio} \qty{8500}{kJ}; \omterm{def:g10:chemistry-of-life:families}{carboidrato} necessário $0.7 \times
8360 \approx \qty{5850}{kJ}$: sim, as reservas bastam sem beber.

\textbf{20.} Sem preparação, o \omterm{def:g10:exercise-and-energy:glycogen}{glicogênio} de Nádia acaba por volta do
quilômetro 34; a carga de \omterm{def:g10:chemistry-of-life:families}{carboidrato} antes da prova, beber \omterm{def:g10:chemistry-of-life:families}{carboidrato}
durante ela e um treino que eleve a parcela de gordura queimada
empurram, cada um, o muro para além da linha de chegada.
\end{solution}
